\documentclass[11pt]{article}
\usepackage[margin=1in]{geometry}
\usepackage{amsmath,amssymb,amsthm}
\usepackage[hidelinks]{hyperref}
\newtheorem{theorem}{Theorem}
\newtheorem{lemma}{Lemma}
\title{Strict convexity of the pendulum time map\\{\large MathDB 343384}}
\author{Principia Math In Progress}
\date{12 September 2026}
\begin{document}
\maketitle

\section*{Task 1. Resolve the time-map convexity conjecture}
For \(0<\phi<\pi/2\) and \(0<z<2\), put
\[
 \beta(z)=\tfrac12\arccos(1-z),\qquad
\]
\[
 T_B(z,\phi)=
 \int_0^\phi\frac{dx}{\sqrt{z-\cos(2\phi)+\cos(2x)}}
 +\int_0^{\beta(z)}
       \frac{dx}{\sqrt{z-\cos(2\phi)+\cos(2x)}}.
\]
Conjecture 1 of \cite{original}, retained as Conjecture 1 in the authors'
postpublication manuscript of \cite{later}, asks whether this is convex in
\(z\) for every fixed \(\phi\).

\begin{theorem}
The conjecture holds. In fact, throughout the stated parameter rectangle,
\[
 \boxed{\displaystyle
 \frac{\partial^2 T_B}{\partial z^2}(z,\phi)
       \geq \frac{\sin\phi}{16z^{3/2}}>0.}
\]
\end{theorem}

\begin{proof}
Fix \(w=1-\cos(2\phi)=2\sin^2\phi\in(0,2)\), and write \(s=z+w\).
In both integrals use \(v=2\sin^2x\). Since
\(dx=dv/(2\sqrt{v(2-v)})\), this gives
\begin{equation}\label{transform}
 T_B(z,\phi)=\frac12
 \left(\int_0^w+\int_0^z\right)
 \frac{dv}{\sqrt{v(2-v)(s-v)}}.
\end{equation}
For each integer \(n\geq0\), define
\[
 F_n(z,w)=
 \int_0^w\frac{v^{n-1/2}}{\sqrt{z+w-v}}\,dv+
 \int_0^z\frac{v^{n-1/2}}{\sqrt{z+w-v}}\,dv.
\]
The binomial series and its nonnegative terms yield
\begin{equation}\label{series}
 T_B(z,\phi)=\frac1{2\sqrt2}
       \sum_{n=0}^{\infty}\frac{\binom{2n}{n}}{8^n}F_n(z,w).
\end{equation}
This passage follows either by monotone convergence or by uniform
convergence of the binomial series on
\(0\leq v\leq\max\{z,w\}<2\).

We prove that \(F_0\) is constant and every \(F_n\), \(n\geq1\), is
convex as a function of \(z>0\) with \(w>0\) fixed.
Indeed,
\[
 F_0(z,w)=2\arcsin\sqrt{\frac{w}{s}}
             +2\arcsin\sqrt{\frac{z}{s}}=\pi.
\]
For \(n\geq1\), the first integral in \(F_n\) has positive second
derivative
\[
 \frac34\int_0^w
       v^{n-1/2}(z+w-v)^{-5/2}\,dv>0.
\]
Differentiation under this fixed integral is justified locally by
\(z+w-v\geq z>0\).
In its second integral substitute \(v=zt\), obtaining
\[
 \int_0^1 t^{n-1/2}
       \frac{z^{n+1/2}}{\sqrt{w+(1-t)z}}\,dt.
\]
For \(b\geq0\), direct differentiation gives
\begin{align}
 \frac{d^2}{dz^2}\frac{z^{n+1/2}}{\sqrt{w+bz}}
 &=\frac{z^{n-3/2}}{(w+bz)^{5/2}}P_n(w,bz),\label{kernel}\\
 P_n(w,u)&=
 (n^2-\tfrac14)w^2+(2n+1)(n-1)wu+n(n-1)u^2.\notag
\end{align}
All coefficients in this polynomial are nonnegative for integer
\(n\geq1\), and its first term is positive. Thus the second integral
is also convex. Its local differentiability follows from \(w>0\) and
the integrable weight \(t^{n-1/2}\).

Consequently the series \eqref{series} is a finite-valued pointwise
limit of convex functions. Convexity passes to this limit by passing
to the limit in the defining convexity inequality.
To get the advertised quantitative conclusion, elementary integration gives
\[
 F_1(z,w)=\frac{\pi}{2}(z+w)-2\sqrt{zw},
 \qquad
 \partial_z^2F_1(z,w)=\frac{\sqrt w}{2z^{3/2}}.
\]
The function
\[
 T_B(z,\phi)-\frac1{8\sqrt2}F_1(z,w)
  =\frac1{2\sqrt2}
       \left(\pi+\sum_{n=2}^{\infty}
                  \frac{\binom{2n}{n}}{8^n}F_n(z,w)\right)
\]
is convex by the same reasoning. It is \(C^2\): the original two
integrals have smooth endpoints and strictly positive radicands on
every compact subset of the parameter rectangle. Its second derivative
is therefore nonnegative. Hence
\[
 \partial_z^2T_B(z,\phi)\geq
       \frac{\sqrt w}{16\sqrt2\,z^{3/2}}
       =\frac{\sin\phi}{16z^{3/2}},
\]
as required.
\end{proof}

\paragraph{Consequence.}
The derivative is strictly increasing. The source proves
\(\partial_zT_B\to-\infty\) as \(z\downarrow0\).
As \(z\uparrow2\), differentiation of the original formula gives
\[
 2\partial_zT_B=
 -\left(\int_0^\phi+\int_0^{\beta(z)}\right)
       (z-\cos2\phi+\cos2x)^{-3/2}\,dx
 +\frac1{\sqrt{z(2-z)w}}\longrightarrow+\infty:
\]
the integrals remain bounded for fixed \(w>0\).
Thus there is exactly one stationary point, a nondegenerate global
minimum. This establishes the single-minimum consequence for the
type \(B\) branch as well.

\paragraph{Verification and status.}
The accompanying verifier checks the algebraic identity
\eqref{kernel}, the first coefficient, the changes of variables, and
the series against direct quadrature. These numerical checks are
supplementary; the proof above covers the entire open rectangle.
The separate literature audit records the live open status, empty
claim fields, current source checks, and the scope of the claim search.

\section*{Task 2. Locate the unique minimum}
Determine asymptotic expansions for its location \(z_*(\phi)\) and
value \(T_B(z_*(\phi),\phi)\) as \(\phi\downarrow0\) and
\(\phi\uparrow\pi/2\), with explicit error estimates.

\section*{Task 3. Resolve the branches with additional turns}
Use the convex expansion together with the complete-period contribution
to settle uniqueness and nondegeneracy of the minima on the higher
turning branches described in \cite{later}, on their full domains.

\section*{Task 4. Extend the convexity mechanism}
Characterize potential functions whose analogous change of variables
produces a positive power-series weight, and establish a time-map
convexity theorem for that class.

\section*{Task 5. Separate the boundary parameters}
Allow the prescribed terminal velocity and initial angle to vary
independently. Determine the parameter region in which the corresponding
mixed-boundary time map is convex and has a unique minimum.

\begin{thebibliography}{9}
\bibitem{original}
F.\ P.\ da Costa, M.\ Grinfeld, J.\ T.\ Pinto, and K.\ Xayxanadasy,
\emph{Bifurcating solutions in a non-homogeneous boundary value problem
for a nonlinear pendulum equation}, arXiv:1907.13009 (2019),
equation (14) and Conjecture 1.
\url{https://arxiv.org/abs/1907.13009}.

\bibitem{later}
F.\ P.\ da Costa, M.\ Grinfeld, N.\ J.\ Mottram, J.\ T.\ Pinto, and
K.\ Xayxanadasy,
\emph{Steady state solutions in a model of a cholesteric liquid crystal
sample}, Afrika Matematika \textbf{32} (2021), 645--672,
doi:10.1007/s13370-020-00851-9.
The authors' postpublication manuscript, including Conjecture 1,
is archived at
\url{https://strathprints.strath.ac.uk/79201/1/da_Costa_etal_AM_2021_Steady_state_solutions_in_a_model.pdf}.
\end{thebibliography}
\end{document}
